\documentclass[11pt]{article}
\usepackage[utf8]{inputenc}
\usepackage[spanish, es-tabla]{babel}
\usepackage{amsmath, amssymb}
\usepackage{booktabs}
\usepackage{enumitem}
\usepackage[margin=2.5cm]{geometry}

\setlength{\parindent}{0pt}
\setlength{\parskip}{0.65em}

\title{Econometría ICS-294\\Guía de ejercicios: Clase 10}
\author{Francisco Alfaro Medina}
\date{}

\begin{document}
\maketitle

\section*{Objetivo}

Practicar inferencia sobre un parámetro individual: hipótesis nula y alternativa, estadístico \(t\), regla de decisión, intervalos de confianza, valor-p y tests de una cola.

\section*{Ejercicios}

\begin{enumerate}[label=\textbf{Ejercicio \arabic*.}, leftmargin=*]

\item \textbf{Hipótesis bilaterales.}

Para cada pregunta, escriba \(H_0\) y \(H_1\).

\begin{enumerate}[label=\alph*)]
    \item ¿La educación tiene efecto sobre salario?
    \item ¿El coeficiente de experiencia es distinto de 2?
    \item ¿El efecto de una política pública es cero?
\end{enumerate}

\item \textbf{Estadístico \(t\).}

En una regresión se obtiene \(\hat{\beta}_1=4.5\), \(se(\hat{\beta}_1)=1.5\). Se quiere probar:
\[
H_0:\beta_1=0
\qquad\text{vs}\qquad
H_1:\beta_1\neq 0.
\]

\begin{enumerate}[label=\alph*)]
    \item Calcule el estadístico \(t\).
    \item Compare con el valor crítico aproximado 1.96.
    \item ¿Rechaza \(H_0\) al 5\%?
\end{enumerate}

\item \textbf{Valor nulo distinto de cero.}

Suponga que \(\hat{\beta}_1=0.12\), \(se(\hat{\beta}_1)=0.04\). Testee:
\[
H_0:\beta_1=0.05
\qquad\text{vs}\qquad
H_1:\beta_1\neq 0.05.
\]

\begin{enumerate}[label=\alph*)]
    \item Calcule el estadístico \(t\).
    \item ¿Rechaza usando valor crítico 1.96?
    \item Interprete el resultado.
\end{enumerate}

\item \textbf{Intervalo de confianza.}

Con \(\hat{\beta}_1=2.4\), \(se(\hat{\beta}_1)=0.6\), construya un intervalo de confianza al 95\% usando 1.96.

\begin{enumerate}[label=\alph*)]
    \item Calcule el límite inferior.
    \item Calcule el límite superior.
    \item ¿El intervalo contiene cero?
    \item ¿Qué implica para un test bilateral al 5\%?
\end{enumerate}

\item \textbf{Valor-p.}

Explique en palabras qué mide el valor-p. Luego ordene los siguientes resultados desde mayor a menor evidencia contra \(H_0\):

\[
p=0.20,\qquad p=0.08,\qquad p=0.03,\qquad p=0.001.
\]

\item \textbf{Test de una cola.}

Se quiere evaluar si la experiencia tiene un efecto positivo:
\[
H_0:\beta_{exper}=0
\qquad\text{vs}\qquad
H_1:\beta_{exper}>0.
\]

Si \(\hat{\beta}_{exper}=0.03\) y \(se(\hat{\beta}_{exper})=0.015\):

\begin{enumerate}[label=\alph*)]
    \item Calcule el estadístico \(t\).
    \item Compare con el valor crítico aproximado 1.645.
    \item ¿Rechaza \(H_0\) al 5\%?
\end{enumerate}

\item \textbf{Errores tipo I y II.}

Explique qué ocurre en cada caso:

\begin{center}
\begin{tabular}{lll}
\toprule
Realidad & Decisión & Resultado \\
\midrule
\(H_0\) verdadera & No rechazar \(H_0\) & \\
\(H_0\) verdadera & Rechazar \(H_0\) & \\
\(H_0\) falsa & No rechazar \(H_0\) & \\
\(H_0\) falsa & Rechazar \(H_0\) & \\
\bottomrule
\end{tabular}
\end{center}

\item \textbf{Interpretación económica.}

Se estima:
\[
\widehat{\log(salario)}=1.2+0.08\,educ+0.01\,exper.
\]
Para educación se obtiene \(se(\hat{\beta}_{educ})=0.02\).

\begin{enumerate}[label=\alph*)]
    \item Calcule el estadístico \(t\) para \(H_0:\beta_{educ}=0\).
    \item ¿Es estadísticamente significativo al 5\%?
    \item Interprete la magnitud del coeficiente.
\end{enumerate}

\item \textbf{Significancia estadística versus económica.}

Discuta si las siguientes afirmaciones son equivalentes:

\begin{enumerate}[label=\alph*)]
    \item Un coeficiente es estadísticamente significativo.
    \item Un coeficiente es económicamente importante.
\end{enumerate}

Proponga un ejemplo donde un coeficiente sea estadísticamente significativo, pero de magnitud pequeña.

\end{enumerate}

\section*{Preguntas de cierre}

\begin{enumerate}[label=\arabic*.]
    \item ¿Qué compara el estadístico \(t\)?
    \item ¿Cuál es la diferencia entre un test bilateral y uno unilateral?
    \item ¿Cómo se relacionan intervalos de confianza y tests de hipótesis?
    \item ¿Por qué un valor-p pequeño entrega evidencia contra \(H_0\)?
\end{enumerate}

\end{document}
